\documentclass[10pt,a4paper]{amsart}
\usepackage[margin=19mm]{geometry}
\usepackage{amsmath,amssymb,mathtools}
\usepackage[scr=rsfs,cal=euler]{mathalfa}
\usepackage{xfrac}
\usepackage[unicode,hidelinks,bookmarks=false]{hyperref}
\newcommand{\ie}{i.e.\ }
\newcommand{\cf}{cf.\ }
\newcommand{\resp}{respectively\ }
\newcommand{\Subsection}{\subsection}
\providecommand{\nobreakdash}{\nobreak\hbox{-}}
\setlength{\emergencystretch}{2em}
\multlinegap0pt
\usepackage{bm}
\usepackage{bbm}
\usepackage{dsfont}
\usepackage{xspace}
\usepackage{cancel}
\usepackage{mathtools}
\usepackage{cleveref}
\usepackage{amsthm}
\makeatletter
\@ifundefined{theorem}{\newtheorem{theorem}{Theorem}}{}
\@ifundefined{lemma}{\newtheorem{lemma}[theorem]{Lemma}}{}
\makeatother
\newcommand{\CC}{\mathbb{C}}
\newcommand{\algmat}[1]{\mathcal{M}(#1)}
\newcommand{\quasigen}{quasi-generic}
\newcommand{\stressmatroid}[2]{\mathcal{S}_{#1}(#2)}
\title[Stress-linked pairs of vertices and the generic stress matro]{Stress-linked pairs of vertices and the generic stress matroid}
\author{D\'aniel Garamv\"olgyi}
\date{Source: arXiv:2308.16851v2 (CC BY 4.0)}
\begin{document}
\maketitle

\noindent\emph{Source and licence.} Stress-linked pairs of vertices and the generic stress matroid. Version arXiv:2308.16851v2 (CC BY 4.0), distributed under the Creative Commons Attribution 4.0 International licence, as recorded in the arXiv licence metadata for this exact version and shown on the abstract page https://arxiv.org/abs/2308.16851v2. Full rights notice: licenses/arxiv-2308.16851-CC-BY-4.0.txt.\par\smallskip

\noindent\textit{Editorial scope.} Counted unit: the Lemma whose source label is \texttt{lemma:projectivetechnical} in the article (source0.tex lines 673--676, source label \texttt{lemma:projectivetechnical}), delivered together with the article's own complete original proof of it (source0.tex lines 679--703, one \texttt{proof} environment of 25 lines). No mathematics is written, repaired or re-derived: statement and proof are the article's own text line for line. Required context retained uncounted: theorem:projectiveimages (lines 664--667). Every definition, display and label of those lines is retained unchanged. Omitted from this package and not redistributed: 1--663, 668--672, 677--678, 704--1473. Nothing inside a retained excerpt is omitted or altered: every retained body is delivered exactly as the article's lines are written, so no omission receipt of any kind accompanies this package. Citations: the retained excerpts carry no citation command at all, so no bibliography entry of the article is transcribed and no reference is redistributed. Reference adaptation: none was needed and none was made; every reference inside the delivered unit resolves inside it, so no unresolved reference or citation is produced. Printed slips kept literal: no printed word, symbol, hypothesis or number of the counted statement or of its proof was altered. Layout and class: the article's own class is \texttt{scrartcl} with packages including inputenc, biblatex, babel, amsmath, amsthm, amssymb, geometry, bm, microtype, mathtools, csquotes, subcaption, enumitem, xcolor; the wrapper is the lane's pinned \texttt{amsart} editorial wrapper at 10pt on A4, which restates the theorem-like environments the retained text needs and adds the article's own macro definitions that the retained text needs (\texttt{\textbackslash CC}, \texttt{\textbackslash algmat}, \texttt{\textbackslash quasigen}, \texttt{\textbackslash stressmatroid}), with extra wrapper packages (bm, bbm, dsfont, xspace, cancel, mathtools, cleveref). The delivered numbering is the unit's own. Assets: no retained span contains a figure environment, a graphics-inclusion command, a PDF-page inclusion, a table or a TikZ picture; no image file is copied into this package, the package declares no asset dependency and no image file is redistributed with it. Licence: version arXiv:2308.16851v2 of this article is distributed under the Creative Commons Attribution 4.0 International licence, as recorded in the arXiv licence metadata for this exact version and shown on the abstract page https://arxiv.org/abs/2308.16851v2. A full rights notice is delivered at licenses/arxiv-2308.16851-CC-BY-4.0.txt. Characters: every retained excerpt is pure ASCII, so the delivered engine, XeLaTeX with its default Latin Modern text font, reports no missing character.
\begin{theorem}\label{theorem:projectiveimages} \emph{(Small forests are independent in the stress matroid.)}\newline
    Let $G = (V,E)$ be a graph, and let $E_0 \subseteq E$ the edge set of a forest $F$ in $G$ on at most $d+1+c$ vertices, where $c$ denotes the number of components of $F$.
    If none of the edges in $E_0$ are $\mathcal{R}_d$-bridges in $G$, then $E_0$ is independent in $\stressmatroid{d}{G}$.
\end{theorem}

\begin{lemma}\label{lemma:projectivetechnical}
    Let $G = (V,E)$ be a graph, let $(G,p)$ be a generic realization in $\CC^d$, and let $\omega$ be a generic $d$-stress with $\omega \in S(G,p)$. Consider the set $P(G,p,\omega) \subseteq (\CC^{d})^V \times \CC^E $ defined as above. %
    We have $\pi_2(P(G,p,\omega)) \subseteq M_{d,G}^*$, where $\pi_2 : (\CC^{d})^V \times \CC^E \rightarrow \CC^{E}$ denotes the projection onto the second factor.
\end{lemma}

\begin{proof}[Proof of \cref{theorem:projectiveimages}]
    Let $\omega$ be a generic $d$-stress of $G$, and let $(G,p)$ be a generic realization in $\CC^d$ with $\omega \in S(G,p)$. Note that $\omega_{uv} \neq 0$ for every $uv \in E_0$. Indeed, since $\omega$ is $M_{d,G}^*$-generic, if it satisfied the equation $x_{uv} = 0$, then every element of $M_{d,G}^*$ would satisfy it. But since $uv$ is not an $\mathcal{R}_d$-bridge of $G$, we know that there exist \quasigen\ $d$-stresses of $G$ that are nonzero on $uv$. 
    
    Let $P(G,p,\omega)$ be defined as above. By \cref{lemma:projectivetechnical}, the projection of $P(G,p,\omega)$ onto $\CC^E$ is contained in $M_{d,G}^*$. Hence, to show that $E_0$ is independent in $\stressmatroid{d}{G} = \algmat{M_{d,G}^*}$, it is enough to show that the projection of $P(G,p,\omega)$ onto $\CC^{E_0}$ is full-dimensional. In other words, it suffices to show that for a nonempty Zariski open set of vectors $\sigma = {(\sigma_e)}_{e \in E_0} \in \CC^{E_0}$, there is a projective image $(q,\omega')$ of $(p,\omega)$ such that $\omega'_e = \sigma_e$ for all $e \in E_0$. For now, let us fix a vector $\sigma \in \CC^{E_0}$ that is everywhere nonzero.
    
    Let us assume for convenience that $c=1$ (i.e., $E_0$ induces a tree) and $|V(E_0)| = d+2$. The proof of the general case is very similar. Let $u_0$ be a leaf vertex in $E_0$, let $v_0$ be its unique neighbor in $E_0$, and set $V_0 = |V(E_0)| - u_0$. Note that $|V_0| = d+1$. 
    
    Recall the notation introduced in the paragraph before \cref{lemma:projectivetechnical}. By genericity, the points $\{p(v), v \in V_0\}$ are affinely independent, and hence the ``lifted'' points $\{\tilde{p}(v), v\in V_0\}$ are linearly independent. This means that we can freely map these points to any collection of $d+1$ points in $\CC^{d+1}$ by using a linear transformation, and in particular, we can freely choose the values $\lambda_v, v \in V_0$ (provided that they are nonzero). Moreover, we can write $\tilde{p}(u_0) = \sum_{v \in V_0}\alpha_v \tilde{p}(v)$ for some nonzero scalars $\alpha_v \in \CC, v \in V_0$. (The fact that $\alpha_v \neq 0$ follows from genericity again.) 
    
    Our goal is to solve the system of equations
    \begin{equation}\label{eq:projective}
        \lambda_u \lambda_v \omega_{uv} = \sigma_{uv}, \hspace{2em} \forall uv \in E_0,
    \end{equation}
    where, by the preceding discussion, $\lambda_v, v \in V_0$ are variables whose value we can freely choose, while by linearity we have $\lambda_{u_0} = \sum_{v \in V_0} \alpha_v \lambda_v$.

    Let us write $\lambda_{v_0} = c$, where $c \in \CC$ is a nonzero scalar whose value we will fix later. For each vertex $v \in V_0 - v_0$, there is a unique path $v_0, v_1, \ldots, v_k = v$ in $F$. We now define $\lambda_v$ by
    \[\lambda_v = 
    \begin{cases}
    \frac{\sigma_{v_0v_1}}{\omega_{v_0v_1}} \cdot \frac{\omega_{v_1v_2}}{\sigma_{v_1v_2}} \cdot \frac{\sigma_{v_2v_3}}{\omega_{v_2v_3}} \cdots \frac{\sigma_{v_{k-1}v_k}}{\omega_{v_{k-1}v_k}} \cdot c  &\text{if } k \text{ is odd,}  \\ \frac{\omega_{v_0v_1}}{\sigma_{v_0v_1}} \cdot \frac{\sigma_{v_1v_2}}{\omega_{v_1v_2}} \cdot \frac{\omega_{v_2v_3}}{\sigma_{v_2v_3}} \cdots \frac{\sigma_{v_{k-1}v_k}}{\omega_{v_{k-1}v_k}} \cdot c^{-1} &\text{if } k \text{ is even.}
    \end{cases}\]

    It is easy to verify that with this definition, \cref{eq:projective} is satisfied for every $uv \in E_0 - u_0v_0$. Moreover, we have \[\lambda_{u_0} = \sum_{v \in V_0}\alpha_v \lambda_v = A \cdot c + B \cdot c^{-1}\] for some scalars $A,B \in \CC$, and hence \[\lambda_{u_0}\lambda_{v_0}\omega_{uv} = (Ac^2 + B) \omega_{uv}.\] If $A \neq 0$, then the right-hand side is surjective as a function of $c$, and hence with a suitable choice of $c$ it can be made to equal $\sigma_{u_0v_0}$. 
    
    Thus it only remains to argue that $A \neq 0$ for a nonempty Zariski open set of vectors $\sigma \in \CC^{E_0}$. After clearing denominators, the condition $A = 0$ leads to a polynomial equation $\Phi(\omega,\sigma) = 0$ involving $\omega$ and $\sigma$. Thus $A \neq 0$ holds on a Zariski open set, and we only need to argue that this set is nonempty, or equivalently, that the polynomial $\Phi$ is not identically zero. This can be seen by looking at a vertex $v \in V_0 - v_0$ whose distance from $v_0$ is odd and maximal among such vertices. Let $e$ be the first edge on the path from $v$ to $v_0$ in $F$. If we look at the expression of $\Phi(\omega,\sigma)$ obtained by clearing denominators in $A$, $\sigma_e$ only appears in a single term. Hence the corresponding variable has positive degree in $\Phi$, and thus $\Phi \neq 0$.
\end{proof}

\end{document}
